\section{Concepto de campo eléctrico}

La ley de Coulomb permite calcular la fuerza que dos cargas ejercen entre sí.
Sin embargo, existe una forma más general de describir la interacción eléctrica.

Consideremos una carga \(Q\) situada en una región del espacio. Si colocamos una
segunda carga \(q\) cerca de ella, aparece una fuerza eléctrica.

Podríamos describir esta situación diciendo simplemente que \(Q\) ejerce una
fuerza a distancia sobre \(q\). No obstante, resulta más útil interpretar que la
presencia de \(Q\) establece una propiedad física en el espacio que la rodea.

A esta propiedad la llamamos \textbf{campo eléctrico}.

\begin{idea}

La carga fuente crea un campo eléctrico en el espacio que la rodea.

Si posteriormente colocamos otra carga en ese campo, la carga experimentará
una fuerza eléctrica.

Por tanto, conviene distinguir entre

\[
\boxed{\text{fuente del campo}}
\]

y

\[
\boxed{\text{carga que experimenta la fuerza}}.
\]

\end{idea}

El concepto de campo permite separar dos problemas:

\begin{enumerate}
    \item determinar qué campo produce una distribución de cargas;
    \item determinar qué fuerza experimenta una carga colocada en ese campo.
\end{enumerate}

Esta separación será muy útil cuando existan varias cargas o distribuciones
continuas de carga.


\section{Definición de campo eléctrico}

Supongamos que queremos conocer el campo eléctrico en un punto \(P\).

Imaginamos que colocamos en \(P\) una pequeña carga positiva \(q_0\), denominada
\textbf{carga de prueba}.

Si sobre ella actúa una fuerza eléctrica \(\vec F_0\), definimos el campo eléctrico
en ese punto como la fuerza eléctrica por unidad de carga de prueba.

\begin{definicion}[Campo eléctrico]

El \textbf{campo eléctrico} en un punto del espacio es una magnitud vectorial
definida como la fuerza eléctrica que experimentaría una carga de prueba positiva
por unidad de carga.

\end{definicion}

Matemáticamente,

\begin{equation}
\boxed{\vec E=\frac{\vec F_0}{q_0}}.
\label{eq:definicion_campo_electrico}
\end{equation}

De manera más rigurosa, la carga de prueba debe ser suficientemente pequeña para
no modificar apreciablemente la distribución de cargas que genera el campo.

Esta idea puede expresarse como

\begin{equation}
\boxed{\vec E=\lim_{q_0\to0}\frac{\vec F_0}{q_0}}.
\end{equation}

\begin{idea}

La carga de prueba no define la existencia del campo.

El campo es producido por las \textbf{cargas fuente} y existe en el punto aunque
no coloquemos allí ninguna carga de prueba.

La carga de prueba es solamente un recurso conceptual para medir o definir el
campo.

\end{idea}


\subsection{Carga fuente y carga de prueba}

Conviene distinguir claramente las dos clases de carga que aparecen en la
definición.

La \textbf{carga fuente} es la carga que produce el campo eléctrico.

La \textbf{carga de prueba} es una carga imaginaria, positiva y suficientemente
pequeña, utilizada para determinar la dirección y la magnitud del campo.

Por ejemplo, si una carga fuente \(Q\) produce un campo en el punto \(P\), el campo

\[
\vec E(P)
\]

depende de \(Q\), de la posición de \(P\) y de las demás cargas fuente existentes.

No depende de la magnitud de la carga de prueba utilizada para investigarlo.

\begin{errorfrecuente}

No interpretar

\[
\vec E=\frac{\vec F_0}{q_0}
\]

como si una carga de prueba grande produjera un campo más pequeño.

Si aumentamos \(q_0\) sin alterar las cargas fuente, también aumenta
proporcionalmente la fuerza:

\[
\vec F_0=q_0\vec E.
\]

Por ello, el cociente permanece igual mientras la carga de prueba no perturbe la
distribución que genera el campo.

\end{errorfrecuente}


\section{El campo eléctrico es una magnitud vectorial}

Como la fuerza eléctrica es vectorial y la carga \(q_0\) es escalar,

\[
\vec E=\frac{\vec F_0}{q_0}
\]

también es una magnitud vectorial.

Para especificar un campo eléctrico en un punto debemos conocer:

\[
\boxed{\text{magnitud, dirección y sentido}}.
\]

Por convención, la dirección del campo eléctrico es la dirección de la fuerza que
experimentaría una \textbf{carga de prueba positiva} colocada en ese punto.

Esta convención es fundamental.

Si una carga positiva experimentaría una fuerza hacia la derecha, entonces el campo
apunta hacia la derecha.

Si experimentaría una fuerza hacia arriba, el campo apunta hacia arriba.


\subsection{Unidad del campo eléctrico}

De

\[
E=\frac{F}{q}
\]

se obtiene

\[
[E]=\frac{\si{\newton}}{\si{\coulomb}}.
\]

Por tanto, en el Sistema Internacional,

\begin{equation}
\boxed{[E]=\si{\newton\per\coulomb}}.
\end{equation}


\section{Fuerza sobre una carga colocada en un campo}

De la definición

\[
\vec E=\frac{\vec F}{q}
\]

podemos despejar la fuerza:

\begin{equation}
\boxed{\vec F=q\vec E}.
\label{eq:fuerza_campo_electrico}
\end{equation}

Esta relación permite calcular la fuerza sobre cualquier carga \(q\) colocada en un
campo eléctrico conocido.


\subsection{Carga positiva}

Si

\[
q>0,
\]

entonces

\[
\vec F=q\vec E
\]

tiene la misma dirección y el mismo sentido que \(\vec E\).

Por tanto,

\[
\boxed{q>0\Longrightarrow\vec F\text{ en el mismo sentido que }\vec E}.
\]


\subsection{Carga negativa}

Si

\[
q<0,
\]

el factor que multiplica a \(\vec E\) es negativo.

Por tanto, la fuerza tiene sentido opuesto al campo:

\[
\boxed{q<0\Longrightarrow\vec F\text{ en sentido opuesto a }\vec E}.
\]

\begin{idea}

El campo eléctrico no cambia de dirección porque coloquemos una carga negativa.

El campo está determinado por las cargas fuente.

Lo que cambia es la dirección de la fuerza que experimenta la carga colocada en ese
campo.

\end{idea}


\subsection{Ejemplo: fuerza sobre una carga en un campo}

\begin{ejemplo}

En cierto punto existe un campo eléctrico uniforme

\[
\vec E=(3.0\times10^4)\,\hat{\mathbf i}\,\si{\newton\per\coulomb}.
\]

Se coloca en ese punto una carga

\[
q=-\SI{2.0}{\nano\coulomb}.
\]

Determine la fuerza eléctrica sobre la carga.

\end{ejemplo}

Convertimos la carga:

\[
q=-\SI{2.0e-9}{\coulomb}.
\]

Aplicamos

\[
\vec F=q\vec E.
\]

Entonces,

\[
\vec F=(-2.0\times10^{-9})(3.0\times10^4)\hat{\mathbf i}\,\si{\newton}.
\]

Por tanto,

\[
\boxed{\vec F=-6.0\times10^{-5}\,\hat{\mathbf i}\,\si{\newton}}.
\]

La magnitud es

\[
\boxed{F=\SI{6.0e-5}{\newton}}.
\]

El signo negativo de la componente indica que la fuerza apunta hacia \(-x\).

Esto es coherente con el hecho de que la carga es negativa: la fuerza tiene sentido
opuesto al campo.


\section{Campo eléctrico producido por una carga puntual}

Consideremos una carga fuente puntual \(Q\).

Queremos determinar el campo eléctrico que produce en un punto \(P\) situado a una
distancia \(r\).

Imaginemos una pequeña carga de prueba positiva \(q_0\) colocada en \(P\).

La magnitud de la fuerza que \(Q\) ejerce sobre \(q_0\) es, por la ley de Coulomb,

\[
F_0=k\frac{|Qq_0|}{r^2}.
\]

Como

\[
E=\frac{F_0}{q_0},
\]

obtenemos

\[
E=\frac{1}{q_0}k\frac{|Qq_0|}{r^2}.
\]

Como \(q_0\) se toma positiva,

\[
\boxed{E=k\frac{|Q|}{r^2}}.
\label{eq:campo_carga_puntual_magnitud}
\]

Esta expresión proporciona la \textbf{magnitud} del campo producido por una carga
puntual.


\subsection{Dirección del campo de una carga puntual}

El campo de una carga puntual es radial.

Si

\[
Q>0,
\]

una carga de prueba positiva sería repelida.

Por tanto,

\[
\boxed{Q>0\Longrightarrow\vec E\text{ apunta alejándose de }Q}.
\]

Si

\[
Q<0,
\]

la carga de prueba positiva sería atraída.

Por tanto,

\[
\boxed{Q<0\Longrightarrow\vec E\text{ apunta hacia }Q}.
\]

Así, una carga positiva produce un campo radial hacia afuera y una carga negativa
produce un campo radial hacia adentro.


\subsection{Dependencia con la distancia}

La magnitud del campo de una carga puntual satisface

\[
E\propto\frac{1}{r^2}.
\]

Por tanto, si la distancia se duplica,

\[
r\to2r,
\]

entonces

\[
\boxed{E\to\frac{E}{4}}.
\]

Si la distancia se triplica,

\[
\boxed{E\to\frac{E}{9}}.
\]

Si la distancia se reduce a la mitad,

\[
\boxed{E\to4E}.
\]

La misma dependencia con \(1/r^2\) que aparece en la fuerza de Coulomb aparece en
el campo producido por una carga puntual.


\subsection{Forma vectorial}

Definamos \(\hat{\mathbf r}\) como un vector unitario dirigido desde la carga fuente
\(Q\) hasta el punto donde calculamos el campo.

Entonces,

\begin{equation}
\boxed{\vec E=k\frac{Q}{r^2}\hat{\mathbf r}}.
\label{eq:campo_carga_puntual_vectorial}
\end{equation}

En esta expresión no usamos valor absoluto.

Si \(Q>0\), el factor es positivo y el campo tiene el mismo sentido que
\(\hat{\mathbf r}\).

Si \(Q<0\), el factor es negativo y el campo apunta en sentido opuesto a
\(\hat{\mathbf r}\), es decir, hacia la carga fuente.


\subsection{Campo mediante vectores de posición}

Supongamos que la carga fuente \(Q\) se encuentra en la posición

\[
\vec r_Q
\]

y queremos determinar el campo en un punto cuya posición es

\[
\vec r.
\]

El vector desde la carga fuente hasta el punto es

\[
\vec R=\vec r-\vec r_Q.
\]

Su magnitud es

\[
R=|\vec r-\vec r_Q|.
\]

El vector unitario es

\[
\hat{\mathbf R}=\frac{\vec r-\vec r_Q}{|\vec r-\vec r_Q|}.
\]

Por tanto,

\begin{equation}
\boxed{
\vec E(\vec r)
=
kQ\frac{\vec r-\vec r_Q}{|\vec r-\vec r_Q|^3}
}.
\label{eq:campo_puntual_posiciones}
\end{equation}

Esta forma es especialmente útil cuando las cargas y los puntos están definidos
mediante coordenadas.


\subsection{Ejemplo: campo de una carga puntual}

\begin{ejemplo}

Una carga puntual

\[
Q=+\SI{4.0}{\micro\coulomb}
\]

produce un campo eléctrico. Determine el campo a una distancia

\[
r=\SI{0.30}{\metre}.
\]

\end{ejemplo}

La magnitud es

\[
E=k\frac{|Q|}{r^2}.
\]

Sustituyendo,

\[
E=(8.99\times10^9)\frac{4.0\times10^{-6}}{(0.30)^2}.
\]

Por tanto,

\[
\boxed{E\approx\SI{4.00e5}{\newton\per\coulomb}}.
\]

Como la carga fuente es positiva, el campo está dirigido radialmente hacia afuera.

\begin{interpretacion}

El resultado no representa una fuerza.

Indica que una carga positiva de \(\SI{1}{\coulomb}\), si pudiera colocarse en ese
punto sin alterar el sistema, experimentaría una fuerza de magnitud
\(\SI{4.00e5}{\newton}\).

Para una carga diferente, la fuerza se obtiene mediante

\[
\vec F=q\vec E.
\]

\end{interpretacion}


\section{Principio de superposición de campos eléctricos}

Si existen varias cargas fuente, cada una produce su propio campo eléctrico en el
punto considerado.

El campo total es la suma vectorial de todos esos campos.

\begin{principio}[Principio de superposición de campos eléctricos]

El campo eléctrico total producido por varias cargas es igual a la suma vectorial
de los campos que cada carga produciría individualmente.

\end{principio}

Para cargas \(q_1,q_2,\ldots,q_n\),

\begin{equation}
\boxed{\vec E=\vec E_1+\vec E_2+\cdots+\vec E_n}.
\end{equation}

De forma compacta,

\begin{equation}
\boxed{\vec E=\sum_i\vec E_i}.
\label{eq:superposicion_campo}
\end{equation}

Cada contribución se calcula mediante

\[
\vec E_i=k\frac{q_i}{r_i^2}\hat{\mathbf r}_i.
\]

Por tanto,

\begin{equation}
\boxed{
\vec E
=
k\sum_i\frac{q_i}{r_i^2}\hat{\mathbf r}_i
}.
\end{equation}


\subsection{Forma general mediante posiciones}

Si la carga \(q_i\) se encuentra en \(\vec r_i\) y queremos conocer el campo en
\(\vec r\),

\[
\boxed{
\vec E(\vec r)
=
k\sum_i
q_i
\frac{\vec r-\vec r_i}{|\vec r-\vec r_i|^3}
}.
\label{eq:campo_superposicion_posiciones}
\]

Esta expresión contiene toda la información geométrica del problema.


\subsection{Superposición en componentes}

En dos dimensiones,

\[
\vec E=E_x\hat{\mathbf i}+E_y\hat{\mathbf j}.
\]

Las componentes totales son

\[
\boxed{E_x=\sum_iE_{i,x}},
\qquad
\boxed{E_y=\sum_iE_{i,y}}.
\]

Si conocemos la magnitud \(E_i\) y el ángulo \(\theta_i\),

\[
E_{i,x}=E_i\cos\theta_i,
\qquad
E_{i,y}=E_i\sin\theta_i.
\]

La magnitud resultante es

\[
\boxed{E=\sqrt{E_x^2+E_y^2}}.
\]

La dirección puede determinarse a partir de las componentes, teniendo siempre en
cuenta el cuadrante.


\subsection{Método para calcular el campo de varias cargas puntuales}

Un procedimiento sistemático es:

\begin{enumerate}
    \item identificar el punto donde se desea calcular el campo;
    \item ubicar todas las cargas fuente;
    \item elegir un sistema de coordenadas;
    \item determinar la distancia desde cada carga hasta el punto;
    \item establecer la dirección de cada campo;
    \item calcular la magnitud de cada contribución;
    \item descomponer los campos en componentes cuando sea necesario;
    \item sumar vectorialmente las contribuciones;
    \item calcular la magnitud y dirección del campo resultante;
    \item comprobar que el sentido obtenido sea físicamente razonable.
\end{enumerate}

\begin{idea}

Para determinar la dirección de \(\vec E_i\), imagine siempre una carga de prueba
\textbf{positiva} en el punto.

Una carga fuente positiva la repelería.

Una carga fuente negativa la atraería.

\end{idea}


\subsection{Ejemplo: campo producido por dos cargas}

\begin{ejemplo}

Una carga \(q_1=+\SI{3.0}{\micro\coulomb}\) se encuentra en el origen y una carga
\(q_2=-\SI{4.0}{\micro\coulomb}\) se encuentra en \((\SI{0.40}{\metre},0)\).

Determine el campo eléctrico en el punto \(P=(0,\SI{0.30}{\metre})\).

\end{ejemplo}

El campo producido por \(q_1\) en \(P\) apunta hacia arriba, porque \(q_1\) es positiva.
El campo producido por \(q_2\) apunta desde \(P\) hacia \(q_2\), porque \(q_2\) es negativa.

\begin{center}
\begin{tikzpicture}[>=Latex,scale=1.1]

\draw[->] (-1.2,0) -- (5.2,0) node[right] {$x$};
\draw[->] (0,-0.8) -- (0,6.0) node[above] {$y$};

\filldraw (0,0) circle (2.2pt) node[below left=3pt] {$q_1$};
\filldraw (4,0) circle (2.2pt) node[below=4pt] {$q_2$};
\filldraw (0,3) circle (2.2pt) node[left=4pt] {$P$};

\draw (4,0.10) -- (4,-0.10) node[below=15pt] {\(\SI{0.40}{\metre}\)};
\draw (0.10,3) -- (-0.10,3) node[left=16pt] {\(\SI{0.30}{\metre}\)};

\draw[dashed] (0,3) -- (4,0);

\draw[->,very thick] (0,3) -- (0,5.2)
node[midway,left=5pt] {$E_1$};

\draw[->,very thick] (0,3) -- (2.4,1.2)
node[midway,below left=2pt] {$E_2$};

\draw[->,thick,dashed] (0,3) -- (1.35,5.5)
node[pos=0.78,right=4pt] {$E$};

\end{tikzpicture}
\end{center}

Para la carga \(q_1\), la distancia hasta \(P\) es

\[
r_1=\SI{0.30}{\metre}.
\]

La magnitud del campo es

\[
E_1=k\frac{|q_1|}{r_1^2}
=(8.99\times10^9)\frac{3.0\times10^{-6}}{(0.30)^2}
\approx\SI{3.00e5}{\newton\per\coulomb}.
\]

Como \(E_1\) apunta verticalmente hacia arriba, sus componentes son

\[
\boxed{E_{1x}=0}, \qquad
\boxed{E_{1y}=+\SI{3.00e5}{\newton\per\coulomb}}.
\]

Para \(q_2\), la distancia desde la carga hasta \(P\) se obtiene mediante

\[
r_2=\sqrt{(0.40)^2+(0.30)^2}=\SI{0.50}{\metre}.
\]

La magnitud del campo producido por \(q_2\) es

\[
E_2=k\frac{|q_2|}{r_2^2}
=(8.99\times10^9)\frac{4.0\times10^{-6}}{(0.50)^2}
\approx\SI{1.44e5}{\newton\per\coulomb}.
\]

La geometría forma un triángulo rectángulo de lados \(0.30\), \(0.40\) y
\(\SI{0.50}{\metre}\), por lo que

\[
\cos\theta=\frac{0.40}{0.50}=0.80,
\qquad
\sin\theta=\frac{0.30}{0.50}=0.60.
\]

Como \(q_2\) es negativa, el campo en \(P\) apunta hacia \(q_2\): su componente
horizontal es positiva y su componente vertical es negativa.

Por tanto,

\[
E_{2x}=E_2\cos\theta=(1.44\times10^5)(0.80)
\approx\SI{1.15e5}{\newton\per\coulomb},
\]

\[
E_{2y}=-E_2\sin\theta=-(1.44\times10^5)(0.60)
\approx-\SI{8.63e4}{\newton\per\coulomb}.
\]

Aplicamos el principio de superposición por componentes.

En el eje \(x\),

\[
E_x=E_{1x}+E_{2x}=0+1.15\times10^5,
\]

por lo que

\[
\boxed{E_x\approx+\SI{1.15e5}{\newton\per\coulomb}}.
\]

En el eje \(y\),

\[
E_y=E_{1y}+E_{2y}=3.00\times10^5-8.63\times10^4,
\]

por lo que

\[
\boxed{E_y\approx+\SI{2.13e5}{\newton\per\coulomb}}.
\]

La magnitud del campo resultante es

\[
E=\sqrt{E_x^2+E_y^2}.
\]

Sustituyendo,

\[
E=\sqrt{(1.15\times10^5)^2+(2.13\times10^5)^2}
\approx\SI{2.42e5}{\newton\per\coulomb}.
\]

Por tanto,

\[
\boxed{E\approx\SI{2.42e5}{\newton\per\coulomb}}.
\]

Como

\[
E_x>0
\qquad\text{y}\qquad
E_y>0,
\]

el campo resultante se encuentra en el primer cuadrante.

Su dirección respecto del eje \(+x\) es

\[
\theta=\tan^{-1}\left(\frac{E_y}{E_x}\right)
=\tan^{-1}\left(\frac{2.13\times10^5}{1.15\times10^5}\right)
\approx61.7^\circ.
\]

Así,

\[
\boxed{\theta\approx61.7^\circ\text{ sobre el eje }+x}.
\]

En conclusión,

\[
\boxed{
E_x\approx+\SI{1.15e5}{\newton\per\coulomb},\qquad
E_y\approx+\SI{2.13e5}{\newton\per\coulomb}
}
\]

y

\[
\boxed{
E\approx\SI{2.42e5}{\newton\per\coulomb},
\qquad
\theta\approx61.7^\circ
}.
\]

\section{Puntos donde el campo eléctrico es nulo}

Puede ocurrir que varias cargas produzcan campos no nulos que se cancelen
vectorialmente en cierto punto.

La condición es

\[
\boxed{\sum_i\vec E_i=\vec 0}.
\]

Esto no significa que las cargas hayan dejado de producir campo.

Significa que sus contribuciones se cancelan exactamente en ese punto.


\subsection{Dos cargas del mismo signo}

Si dos cargas tienen el mismo signo, entre ellas los campos apuntan en sentidos
opuestos.

Por ello, puede existir un punto entre las cargas donde

\[
E_1=E_2.
\]

El punto de cancelación se encuentra más cerca de la carga de menor magnitud.


\subsection{Dos cargas de signos opuestos}

Entre dos cargas de signos opuestos, los campos apuntan en el mismo sentido.

Por tanto, no pueden cancelarse entre ellas en la región comprendida entre las
cargas.

Si existe un punto de campo nulo sobre la línea que las une, deberá encontrarse
fuera del segmento y, en general, del lado de la carga de menor magnitud.

\begin{advertencia}

No basta con igualar magnitudes:

\[
E_1=E_2.
\]

Antes debe comprobarse que los dos vectores tengan sentidos opuestos en la región
analizada.

\end{advertencia}


\section{Distribuciones continuas de carga}

Hasta ahora hemos considerado cargas puntuales.

Sin embargo, en muchos cuerpos macroscópicos la carga está distribuida sobre una
línea, una superficie o un volumen.

Cuando existen enormes cantidades de partículas cargadas muy próximas entre sí,
resulta útil describir la carga mediante una distribución continua.

Dividimos conceptualmente la distribución en elementos infinitesimales de carga

\[
dq.
\]

Cada elemento produce una contribución infinitesimal al campo

\[
d\vec E.
\]

El campo total se obtiene sumando todas las contribuciones.

En el límite continuo, esta suma se convierte en una integral:

\begin{equation}
\boxed{\vec E=\int d\vec E}.
\label{eq:campo_distribucion_continua}
\end{equation}


\subsection{Densidad lineal de carga}

Si la carga se distribuye a lo largo de una línea, definimos la densidad lineal

\begin{equation}
\boxed{\lambda=\frac{dq}{d\ell}}.
\end{equation}

Por tanto,

\begin{equation}
\boxed{dq=\lambda\,d\ell}.
\end{equation}

Su unidad SI es

\[
\boxed{[\lambda]=\si{\coulomb\per\metre}}.
\]

Si la distribución es uniforme sobre una longitud \(L\),

\[
\boxed{\lambda=\frac{Q}{L}}.
\]


\subsection{Densidad superficial de carga}

Si la carga está distribuida sobre una superficie,

\begin{equation}
\boxed{\sigma=\frac{dq}{dA}}.
\end{equation}

Entonces,

\begin{equation}
\boxed{dq=\sigma\,dA}.
\end{equation}

Su unidad es

\[
\boxed{[\sigma]=\si{\coulomb\per\metre\squared}}.
\]

Para una distribución uniforme sobre un área \(A\),

\[
\boxed{\sigma=\frac{Q}{A}}.
\]


\subsection{Densidad volumétrica de carga}

Si la carga ocupa un volumen,

\begin{equation}
\boxed{\rho=\frac{dq}{dV}}.
\end{equation}

Por tanto,

\begin{equation}
\boxed{dq=\rho\,dV}.
\end{equation}

Su unidad es

\[
\boxed{[\rho]=\si{\coulomb\per\metre\cubed}}.
\]

Para una distribución uniforme en un volumen \(V\),

\[
\boxed{\rho=\frac{Q}{V}}.
\]

\begin{advertencia}

Los símbolos

\[
\lambda,\qquad\sigma,\qquad\rho
\]

representan densidades de carga diferentes.

Debe identificarse primero si la distribución es lineal, superficial o
volumétrica antes de escribir \(dq\).

\end{advertencia}


\section{Campo de una distribución continua}

Sea un elemento de carga \(dq\).

La magnitud del campo infinitesimal que produce en un punto situado a distancia
\(R\) es

\[
dE=k\frac{|dq|}{R^2}.
\]

Vectorialmente,

\begin{equation}
\boxed{d\vec E=k\frac{dq}{R^2}\hat{\mathbf R}},
\end{equation}

donde \(\hat{\mathbf R}\) apunta desde el elemento de carga hacia el punto donde
calculamos el campo.

El campo total es

\begin{equation}
\boxed{
\vec E
=
k\int\frac{dq}{R^2}\hat{\mathbf R}
}.
\label{eq:integral_campo_general}
\end{equation}

En términos de posiciones, si \(\vec r'\) representa la posición del elemento
fuente y \(\vec r\) el punto donde calculamos el campo,

\begin{equation}
\boxed{
\vec E(\vec r)
=
k\int
\frac{\vec r-\vec r'}
{|\vec r-\vec r'|^3}\,dq
}.
\label{eq:campo_continuo_vectorial}
\end{equation}


\subsection{Distribución lineal}

Como

\[
dq=\lambda\,d\ell,
\]

tenemos

\begin{equation}
\boxed{
\vec E
=
k\int
\frac{\lambda\,d\ell}{R^2}\hat{\mathbf R}
}.
\end{equation}


\subsection{Distribución superficial}

Como

\[
dq=\sigma\,dA,
\]

obtenemos

\begin{equation}
\boxed{
\vec E
=
k\int
\frac{\sigma\,dA}{R^2}\hat{\mathbf R}
}.
\end{equation}


\subsection{Distribución volumétrica}

Como

\[
dq=\rho\,dV,
\]

resulta

\begin{equation}
\boxed{
\vec E
=
k\int
\frac{\rho\,dV}{R^2}\hat{\mathbf R}
}.
\end{equation}


\section{Simetría en distribuciones continuas}

Antes de integrar es fundamental analizar la geometría.

La simetría puede indicar qué componentes del campo se cancelan.

Por ejemplo, si dos elementos de carga simétricos producen campos

\[
d\vec E_1
\qquad\text{y}\qquad
d\vec E_2,
\]

puede ocurrir que

\[
dE_{1,y}=-dE_{2,y},
\]

de modo que

\[
dE_{1,y}+dE_{2,y}=0.
\]

Mientras tanto, sus componentes en otra dirección pueden sumarse.

\begin{idea}

En problemas con distribuciones continuas no conviene comenzar inmediatamente con
la integral.

Primero debe preguntarse:

\begin{enumerate}
    \item ¿qué simetría tiene la distribución?;
    \item ¿qué componentes se cancelan?;
    \item ¿en qué dirección debe apuntar el campo resultante?;
    \item ¿qué variable conviene utilizar para describir \(dq\)?
\end{enumerate}

La física debe determinar la estructura de la integral antes de realizar el cálculo.

\end{idea}


\section{Ejemplo: campo sobre el eje de un anillo cargado}

\begin{ejemplo}

Un anillo de radio \(a\) posee una carga total \(Q\) distribuida uniformemente.

Determine el campo eléctrico en un punto situado sobre el eje del anillo a una
distancia \(x\) de su centro.

\end{ejemplo}

Consideremos un pequeño elemento de carga \(dq\).

La distancia desde cualquier elemento del anillo hasta el punto es la misma:

\[
R=\sqrt{x^2+a^2}.
\]

El elemento produce un campo de magnitud

\[
dE=k\frac{dq}{R^2}.
\]

Por simetría, a cada elemento del anillo le corresponde otro elemento opuesto cuyas
componentes perpendiculares al eje se cancelan.

Por tanto, solamente sobreviven las componentes paralelas al eje.

Si \(\theta\) es el ángulo entre \(d\vec E\) y el eje,

\[
dE_x=dE\cos\theta.
\]

Geométricamente,

\[
\cos\theta=\frac{x}{R}.
\]

Entonces,

\[
dE_x
=
k\frac{dq}{R^2}\frac{x}{R}
=
k\frac{x\,dq}{R^3}.
\]

Como \(R\) es constante para todos los elementos,

\[
E_x
=
k\frac{x}{R^3}\int dq.
\]

Pero

\[
\int dq=Q.
\]

Por tanto,

\[
E_x=k\frac{Qx}{R^3}.
\]

Sustituyendo

\[
R=\sqrt{x^2+a^2},
\]

obtenemos

\begin{equation}
\boxed{
\vec E
=
k\frac{Qx}{(x^2+a^2)^{3/2}}\hat{\mathbf x}
}.
\label{eq:campo_anillo_eje}
\end{equation}

\begin{interpretacion}

En el centro del anillo,

\[
x=0,
\]

por lo que

\[
\boxed{\vec E=\vec 0}.
\]

Esto también puede deducirse directamente mediante simetría.

Para puntos muy alejados,

\[
x\gg a,
\]

tenemos aproximadamente

\[
(x^2+a^2)^{3/2}\approx x^3.
\]

Entonces,

\[
E\approx k\frac{Q}{x^2}.
\]

A grandes distancias, el anillo se comporta aproximadamente como una carga puntual
de carga total \(Q\).

\end{interpretacion}


\section{Campo uniforme y campo no uniforme}

Un campo eléctrico es \textbf{uniforme} en una región si el vector \(\vec E\) tiene
la misma magnitud, dirección y sentido en todos los puntos de esa región.

Es decir,

\[
\boxed{\vec E=\text{constante}}.
\]

En una representación mediante líneas de campo, un campo uniforme se dibuja mediante
líneas:

\begin{itemize}
    \item paralelas;
    \item igualmente separadas;
    \item orientadas en el mismo sentido.
\end{itemize}

Un campo es \textbf{no uniforme} cuando su magnitud, su dirección o ambas cambian con
la posición.

El campo de una carga puntual es un ejemplo de campo no uniforme porque

\[
E=k\frac{|Q|}{r^2}
\]

depende de la distancia.


\section{Líneas de campo eléctrico}

El campo eléctrico es un vector definido en cada punto del espacio.

Para visualizar una distribución completa de vectores sería necesario dibujar una
gran cantidad de flechas.

Una representación más conveniente consiste en utilizar
\textbf{líneas de campo eléctrico}.

\begin{definicion}[Línea de campo eléctrico]

Una \textbf{línea de campo eléctrico} es una curva imaginaria cuya tangente en cada
punto tiene la misma dirección que el vector campo eléctrico en ese punto.

\end{definicion}

Las líneas no son objetos materiales.

Son una representación gráfica del campo.


\subsection{Dirección de las líneas}

En cualquier punto,

\[
\boxed{
\vec E
\text{ es tangente a la línea de campo}
}.
\]

Las flechas sobre las líneas indican el sentido del campo.


\subsection{Líneas de una carga positiva}

Para una carga puntual positiva, el campo apunta radialmente hacia afuera.

Por tanto, las líneas de campo

\[
\boxed{\text{salen de las cargas positivas}}.
\]


\subsection{Líneas de una carga negativa}

Para una carga puntual negativa, el campo apunta radialmente hacia la carga.

Por tanto,

\[
\boxed{\text{las líneas entran en las cargas negativas}}.
\]


\subsection{Intensidad y densidad de líneas}

La separación entre las líneas permite representar cualitativamente la magnitud del
campo.

Donde las líneas están más próximas,

\[
\boxed{E\text{ es mayor}}.
\]

Donde las líneas están más separadas,

\[
\boxed{E\text{ es menor}}.
\]

Para una carga puntual, las líneas aparecen cada vez más separadas conforme aumenta
la distancia, en concordancia con

\[
E\propto\frac{1}{r^2}.
\]

\begin{advertencia}

El número absoluto de líneas que dibujamos es arbitrario.

Lo importante es conservar una representación coherente en la que una carga de mayor
magnitud tenga proporcionalmente más líneas asociadas.

\end{advertencia}


\section{Reglas para dibujar líneas de campo}

Las líneas de campo electrostático deben cumplir varias propiedades.

\begin{enumerate}
    \item Salen de cargas positivas y terminan en cargas negativas.

    \item Si existe exceso de carga positiva, algunas líneas pueden extenderse hasta
    el infinito.

    \item Si existe exceso de carga negativa, algunas líneas pueden provenir del
    infinito.

    \item La tangente a una línea en cada punto indica la dirección de \(\vec E\).

    \item Una mayor densidad de líneas representa un campo de mayor magnitud.

    \item El número relativo de líneas asociado a una carga debe ser proporcional a
    la magnitud de la carga.

    \item Dos líneas de campo eléctrico no pueden cruzarse.

    \item En electrostática, las líneas no forman curvas cerradas.
\end{enumerate}


\subsection{¿Por qué las líneas no pueden cruzarse?}

Supongamos que dos líneas de campo se cruzaran en un punto.

Cada línea tendría una tangente diferente en ese punto.

Esto implicaría que el campo eléctrico tendría simultáneamente dos direcciones
distintas en el mismo lugar.

Pero en un punto determinado el vector

\[
\vec E
\]

tiene una única dirección.

Por tanto,

\[
\boxed{\text{dos líneas de campo no pueden cruzarse}}.
\]


\section{Líneas de campo de dos cargas}

\subsection{Dos cargas de signos opuestos}

Para dos cargas de igual magnitud y signos opuestos, muchas líneas salen de la carga
positiva y terminan en la negativa.

Esta configuración recibe el nombre de \textbf{dipolo eléctrico}.

Cerca de cada carga las líneas son aproximadamente radiales.

Entre las cargas existe una alta concentración de líneas, indicando que el campo
puede ser intenso en esa región.


\subsection{Dos cargas del mismo signo}

Si las dos cargas son positivas, las líneas salen de ambas.

Las líneas no pueden conectarse de una carga positiva a otra positiva porque el campo
apunta alejándose de ambas.

En una configuración de dos cargas positivas iguales existe un punto medio donde,
por simetría,

\[
\vec E=\vec 0.
\]

Las líneas se curvan alejándose de la región entre las cargas.

Para dos cargas negativas iguales ocurre una configuración análoga, pero las flechas
apuntan hacia las cargas.


\section{Las líneas de campo no son trayectorias}

Una confusión frecuente consiste en interpretar las líneas de campo como caminos que
necesariamente siguen las partículas cargadas.

Esto no es correcto.

Las líneas indican la dirección instantánea del campo eléctrico en cada punto.

La trayectoria de una partícula depende de su:

\begin{itemize}
    \item carga;
    \item masa;
    \item posición inicial;
    \item velocidad inicial;
    \item fuerzas adicionales que puedan actuar.
\end{itemize}

Por ejemplo, una carga negativa experimenta una fuerza en sentido opuesto a
\(\vec E\), por lo que no seguiría las flechas de las líneas de campo.

Incluso una carga positiva con velocidad inicial puede describir una trayectoria que
no coincida con una línea de campo.

\begin{idea}

\[
\boxed{\text{línea de campo}\neq\text{trayectoria de una partícula}}
\]

Una línea de campo representa una propiedad del espacio.

Una trayectoria representa el movimiento de un objeto.

\end{idea}


\section{Campo eléctrico y fuerza eléctrica}

Es importante distinguir los dos conceptos.

La fuerza eléctrica

\[
\vec F
\]

describe la interacción que experimenta una carga específica.

El campo eléctrico

\[
\vec E
\]

describe una propiedad de una región del espacio producida por las cargas fuente.

La relación entre ambos es

\[
\boxed{\vec F=q\vec E}.
\]

Esto significa que en un mismo punto pueden colocarse diferentes cargas y obtenerse
diferentes fuerzas, aunque el campo sea el mismo.

Por ejemplo, si

\[
q_2=2q_1,
\]

entonces, dentro del mismo campo,

\[
\vec F_2=2\vec F_1.
\]

Pero

\[
\vec E
\]

no cambia por ese hecho, siempre que las cargas introducidas no alteren
significativamente las cargas fuente.


\section{Errores y confusiones frecuentes}

\begin{errorfrecuente}

\textbf{Confundir campo eléctrico con fuerza eléctrica.}

El campo se mide en

\[
\si{\newton\per\coulomb},
\]

mientras que la fuerza se mide en

\[
\si{\newton}.
\]

Están relacionados mediante

\[
\vec F=q\vec E.
\]

\end{errorfrecuente}

\begin{errorfrecuente}

\textbf{Pensar que el campo depende de la carga de prueba.}

El campo está determinado por las cargas fuente.

La carga de prueba solamente permite definir o medir el campo.

\end{errorfrecuente}

\begin{errorfrecuente}

\textbf{Determinar la dirección del campo utilizando una carga de prueba negativa.}

Por convenio, la dirección de \(\vec E\) es la dirección de la fuerza que
experimentaría una carga de prueba \textbf{positiva}.

\end{errorfrecuente}

\begin{errorfrecuente}

\textbf{Pensar que una carga negativa produce un campo de magnitud negativa.}

La magnitud

\[
E=k\frac{|Q|}{r^2}
\]

es siempre no negativa.

El signo de \(Q\) determina el sentido del vector campo.

\end{errorfrecuente}

\begin{errorfrecuente}

\textbf{Pensar que una carga negativa colocada en un campo se mueve necesariamente
en la dirección del campo.}

La fuerza sobre una carga negativa apunta en sentido opuesto:

\[
\vec F=q\vec E,
\qquad q<0.
\]

\end{errorfrecuente}

\begin{errorfrecuente}

\textbf{Sumar magnitudes de campos que tienen diferentes direcciones.}

La superposición es vectorial:

\[
\vec E=\sum_i\vec E_i.
\]

\end{errorfrecuente}

\begin{errorfrecuente}

\textbf{Asignar el signo de una componente según el signo de la carga fuente.}

Los signos de \(E_x\) y \(E_y\) dependen de la dirección del vector respecto de los
ejes elegidos.

\end{errorfrecuente}

\begin{errorfrecuente}

\textbf{Olvidar que la distancia en el campo de una carga puntual aparece al
cuadrado.}

\[
E=k\frac{|Q|}{r^2}.
\]

Duplicar \(r\) reduce \(E\) a la cuarta parte.

\end{errorfrecuente}

\begin{errorfrecuente}

\textbf{Igualar campos para buscar un punto de cancelación sin analizar primero sus
direcciones.}

Para que

\[
\vec E_1+\vec E_2=\vec 0,
\]

no basta con que

\[
E_1=E_2.
\]

También deben tener sentidos opuestos.

\end{errorfrecuente}

\begin{errorfrecuente}

\textbf{Pensar que las líneas de campo son objetos físicos.}

Las líneas son una representación gráfica del campo y el número exacto de líneas
dibujadas es una elección del esquema.

\end{errorfrecuente}

\begin{errorfrecuente}

\textbf{Pensar que las líneas de campo son siempre trayectorias de partículas.}

Las líneas representan la dirección del campo, no necesariamente el movimiento de una
partícula.

\end{errorfrecuente}


\section{Preguntas para comprobar comprensión}

\begin{enumerate}
    \item ¿Qué diferencia existe entre una carga fuente y una carga de prueba?

    \item ¿Puede existir un campo eléctrico en un punto donde no haya ninguna carga?

    \item ¿Por qué la carga de prueba debe ser suficientemente pequeña?

    \item ¿Qué representa físicamente la relación
    \[
    \vec E=\frac{\vec F}{q_0}?
    \]

    \item ¿Cuál es la unidad SI del campo eléctrico?

    \item ¿Qué dirección se utiliza para definir el vector campo eléctrico?

    \item Si colocamos una carga positiva en un campo, ¿en qué dirección actúa la
    fuerza?

    \item Si colocamos una carga negativa, ¿en qué dirección actúa la fuerza?

    \item ¿Una carga negativa produce un campo eléctrico negativo?

    \item ¿En qué dirección apunta el campo alrededor de una carga puntual positiva?

    \item ¿Y alrededor de una carga puntual negativa?

    \item Si duplicamos la distancia a una carga puntual, ¿cómo cambia el campo?

    \item ¿Qué principio utilizamos para determinar el campo de varias cargas?

    \item ¿Por qué los campos deben sumarse vectorialmente?

    \item ¿Puede el campo total ser cero aunque cada carga produzca un campo distinto
    de cero?

    \item ¿Qué representan \(\lambda\), \(\sigma\) y \(\rho\)?

    \item ¿Por qué usamos una integral para calcular el campo de una distribución
    continua?

    \item ¿Qué información puede obtenerse de la simetría antes de integrar?

    \item ¿Qué representa la tangente a una línea de campo?

    \item ¿Qué significa que las líneas estén muy próximas entre sí?

    \item ¿Por qué las líneas de campo no pueden cruzarse?

    \item ¿Por qué las líneas salen de cargas positivas y entran en cargas negativas?

    \item ¿Una línea de campo representa necesariamente la trayectoria de una
    partícula cargada?

    \item ¿Qué características gráficas tendría un campo eléctrico uniforme?
\end{enumerate}


\begin{resumen}

\begin{itemize}
    \item El campo eléctrico describe la influencia eléctrica producida por una
    distribución de cargas en una región del espacio.

    \item Se define mediante
    \[
    \boxed{\vec E=\frac{\vec F_0}{q_0}},
    \]
    donde \(q_0\) es una carga de prueba positiva suficientemente pequeña.

    \item El campo es una magnitud vectorial y su unidad SI es
    \[
    \boxed{\si{\newton\per\coulomb}}.
    \]

    \item La fuerza sobre una carga colocada en un campo es
    \[
    \boxed{\vec F=q\vec E}.
    \]

    \item Para una carga positiva, la fuerza tiene el mismo sentido que el campo; para
    una carga negativa, tiene sentido opuesto.

    \item El campo producido por una carga puntual tiene magnitud
    \[
    \boxed{E=k\frac{|Q|}{r^2}}.
    \]

    \item Vectorialmente,
    \[
    \boxed{\vec E=k\frac{Q}{r^2}\hat{\mathbf r}}.
    \]

    \item El campo apunta alejándose de una carga positiva y hacia una carga
    negativa.

    \item El campo de varias cargas se obtiene mediante superposición:
    \[
    \boxed{\vec E=\sum_i\vec E_i}.
    \]

    \item En términos de las posiciones de cargas puntuales,
    \[
    \boxed{
    \vec E(\vec r)
    =
    k\sum_i q_i
    \frac{\vec r-\vec r_i}{|\vec r-\vec r_i|^3}
    }.
    \]

    \item Para distribuciones continuas,
    \[
    \boxed{\vec E=\int d\vec E}.
    \]

    \item Las densidades de carga son
    \[
    \boxed{\lambda=\frac{dq}{d\ell}},
    \qquad
    \boxed{\sigma=\frac{dq}{dA}},
    \qquad
    \boxed{\rho=\frac{dq}{dV}}.
    \]

    \item La expresión general para una distribución continua es
    \[
    \boxed{
    \vec E(\vec r)
    =
    k\int
    \frac{\vec r-\vec r'}
    {|\vec r-\vec r'|^3}\,dq
    }.
    \]

    \item Las líneas de campo son tangentes al vector \(\vec E\) en cada punto.

    \item Las líneas salen de cargas positivas y terminan en cargas negativas o en el
    infinito.

    \item Una mayor densidad de líneas representa un campo más intenso.

    \item Dos líneas de campo no pueden cruzarse.

    \item Las líneas de campo no deben confundirse con trayectorias de partículas.

    \item Un campo uniforme posee la misma magnitud y dirección en todos los puntos
    de la región considerada.
\end{itemize}

\end{resumen}